\section{The integral model and exact boundary conditions}\label{sec:model}

We specify the integral model once and record how to select exact
boundary classes on its complex quotient. In particular, we distinguish
semisimplification of a representation from diagonalization of its
individual matrices.

\subsection{Representation schemes}

For a finitely presented group $\Gamma$ and an integer $r\geq1$, let
$\mathscr R_r(\Gamma)$ be the affine $\Z$-scheme representing
\[
 S\longmapsto\Hom(\Gamma,\SL_r(S)).
\]
Concretely, choose a matrix for each generator and impose the matrix
relations of a finite presentation. Simultaneous conjugation defines
an action of $\SL_{r,\Z}$. With
$A=\Z[\mathscr R_r(\Gamma)]$, put
\begin{equation}\label{eq:integral-model}
 M_r(\Gamma)=\Spec\bigl(A^{\SL_{r,\Z}}\bigr).
\end{equation}
This definition is independent of the chosen presentation, by the
representing property. The invariant ring is finitely generated over
$\Z$ by finite generation for Chevalley-group invariants
\cite[Theorem~3]{FranjouVanderKallen}.
Here the base $\Z$ is Noetherian, $A$ is a finite-type commutative
algebra, and $\SL_{r,\Z}$ is a split Chevalley group with its
rational conjugation action.  This input does not require $A$ to
be reduced or flat over $\Z$.

Flat base change commutes with these invariants: they are the kernel of
the difference between the coaction and the map $a\mapsto a\otimes1$,
and tensoring with a flat algebra preserves that kernel. In particular,
the complex fiber of \eqref{eq:integral-model} is the usual affine
character quotient. Its complex points correspond to the unique closed orbits
in the fibers of the quotient~\cite[pp.~184--185]{MumfordSuominen}.
The representation scheme is closed in the space of generator
tuples: inverses in $\SL_r$ are polynomial adjugates, so the
relators are polynomial equations. Kraft's matrix-tuple closed-orbit
criterion, recalled in \cite[Section~3.4]{BateMartinRoehrleTange},
therefore identifies these orbits with semisimple representations.
Conjugation by $\GL_r(\C)$ gives the same orbits as conjugation by
$\SL_r(\C)$, since any invertible matrix can be multiplied by a scalar
to have determinant one.

For the curve $X$, we use $\Gamma=\pi_1(X)$ and write $M_r(X)=M_r(\Gamma)$.
Every representation $\Gamma\to\SL_r(\cO_L)$ gives a point of
$M_r(X)(\cO_L)$ by evaluating invariant functions. Thus it suffices to
construct actual integral matrices. No assertion about generation of
invariants in positive characteristic is needed for this implication.

\subsection{Rank conditions on the quotient}

Write $\mathscr R=\mathscr R_r(\Gamma)_\C$ and let
$q:\mathscr R\to M_r(\Gamma)_\C$ be the quotient map. We use reduced
closed subsets when imposing additional conditions. An invariant closed
subset $D\subset\mathscr R$ has closed image under $q$. Moreover,
\begin{equation}\label{eq:closed-orbit-test}
 q(x)\in q(D)
 \quad\Longleftrightarrow\quad
 \text{the closed orbit in }q^{-1}(q(x))\text{ lies in }D.
\end{equation}
Indeed the closure of every orbit in the fiber contains that closed
orbit, and a closed invariant subset contains the closure of each of
its orbits. These are the standard closed-orbit properties of an affine
quotient by a reductive group in characteristic zero~\cite[pp.~184--185]{MumfordSuominen}.

Fix the characteristic polynomial of each boundary matrix and let
$\lambda$ range over its eigenvalues. A Jordan form is determined by
these eigenvalues and the numbers
\[
 d_{i,\lambda,k}=\rk\bigl((C_i-\lambda I)^k\bigr),
 \qquad 1\leq k\leq r.
\]
Here the expression on the right means the rank for any matrix in
the chosen class $C_i$. Let $D$ be the invariant closed subset defined
by the fixed characteristic polynomials and the inequalities
\[
 \rk\bigl((\rho(\gamma_i)-\lambda I)^k\bigr)
 \leq d_{i,\lambda,k}.
\]
Let $D_{\mathrm{sm}}\subset D$ be the finite union of the closed subsets
on which at least one of these inequalities is strict; indices with
$d_{i,\lambda,k}=0$ contribute an empty subset. The matrices have exactly
the prescribed classes precisely on $D\setminus D_{\mathrm{sm}}$.

\begin{lemma}\label{lem:exact-stratum}
The locus of characters whose semisimple representations have exactly
the prescribed boundary classes is
\[
 Y=q(D)\setminus q(D_{\mathrm{sm}}).
\]
In particular, $Y$ is open in the closed subset $q(D)$ and is defined
over any number field containing the boundary eigenvalues.
\end{lemma}

\begin{proof}
Apply \eqref{eq:closed-orbit-test} first to $D$, then to
$D_{\mathrm{sm}}$. A closed orbit belongs to the first set and not the
second exactly when its boundary matrices have all the prescribed
ranks. These ranks determine their Jordan forms. The defining
equations and rank inequalities have coefficients in the indicated
field $K$.  Write $A_K=A\otimes_{\Z}K$.  If $I_K\subset A_K$ is
the ideal of either closed rank locus, its closed quotient image is defined by the
contracted ideal $I_K\cap A_K^{\SL_r}$.  This is the kernel of
$A_K^{\SL_r}\to A_K/I_K$; flat extension from $K$ to $\C$ preserves
that kernel and the invariant algebra.  The two closed images, and
hence their difference, are therefore defined over $K$.
\end{proof}

The lemma permits us to construct integral representations in closed
rank bounds and impose the exact conditions after passing to characters.
Density will then meet every open part of the desired locus, including
those in individual irreducible components.
